\chapter{El mecanismo Cherenkov}
Las tablas de Step~2 y Step~4 separan la fracción Cherenkov del primer fotón de la población total. El ángulo de emisión y el ángulo crítico se calculan por longitud de onda cuando el runtime es dispersivo. Para un índice $n$, la identidad geométrica usada es
\[
\arccos(\sin\theta_C)=\arcsin(1/n)=\theta_{crit}.
\]
La condición geométrica de atrapamiento del cono se tabula junto con la velocidad axial del borde, sin parámetros libres adicionales.

La caústica angular aparece en las distribuciones de ángulo de salida. La penalización temporal se tabula con la forma $\Delta t=d/v_g\,[\sec\alpha-\sec\alpha_{edge}]$. Las figuras se separan por campaña: las de `step4/` son baseline y las de `step6\_v2/step4/` son v2; todas las captions de este capítulo deben conservar ese prefijo al incorporar las figuras.

La estadística de orden de emisión y los contrastes de multiplicidad están en los CSV de Step~4. La separación Cherenkov/centelleo se conserva como observable tabulado, no como una explicación única de todos los residuos.

\begin{figure}[ht]
\centering
\includegraphics[width=0.82\textwidth]{v2_cherenkov_edge_caustic.pdf}
\caption{Caústica angular de la campaña v2 con óptica dispersiva.}
\end{figure}

\begin{figure}[ht]
\centering
\includegraphics[width=0.82\textwidth]{baseline_cherenkov_guiding_diagnostics.pdf}
\caption{Diagnóstico de guiado Cherenkov de la campaña baseline con índice constante.}
\end{figure}
